\documentclass[12pt]{article}
 
%Russian-specific packages
%--------------------------------------
\usepackage[T2A]{fontenc}
\usepackage[utf8]{inputenc}
\usepackage[english, russian]{babel}
%for search in russian
\usepackage{cmap}
%--------------------------------------

%Math-specific packages
%--------------------------------------
\usepackage{amsmath}
\usepackage{amssymb}

%Format-specific packages
%--------------------------------------
\usepackage[left=2cm,
            right=2cm,
            top=1cm,
            bottom=2cm,
            bindingoffset=0cm]{geometry}

\usepackage{amsthm}
\newtheorem{lemma}{Лемма}
\newtheorem{definition}{Определение}
\newtheorem*{remark}{Замечание}

\def\LCS{
    \mathrm{LCS}
}

\usepackage{tikz}
\usetikzlibrary{tikzmark}
\usetikzlibrary{matrix}
\tikzset{ 
table/.style={
  matrix of nodes,
  row sep=-\pgflinewidth,
  column sep=-\pgflinewidth,
  nodes={draw, rectangle,text width=3em,align=center},
  text depth=1.25ex,
  text height=2.5ex,
  nodes in empty cells
},
row 1/.style={nodes={fill=green!10,text depth=0.4ex,text height=2ex}},
column 1/.style={nodes={fill=green!10}},
}

\usepackage{listings}
\usepackage{xcolor}

%New colors defined below
\definecolor{codegreen}{rgb}{0,0.6,0}
\definecolor{codegray}{rgb}{0.5,0.5,0.5}
\definecolor{codepurple}{rgb}{0.58,0,0.82}
\definecolor{backcolour}{rgb}{0.95,0.95,0.92}

\lstset {
    language=c++,
    backgroundcolor=\color{backcolour},
    commentstyle=\color{codegreen},
    keywordstyle=\color{magenta},
    numberstyle=\tiny\color{codegray},
    stringstyle=\color{codepurple},
    basicstyle=\ttfamily\footnotesize,
    captionpos=b,
    keepspaces=true,
    numbers=left,
    numbersep=5pt,
    tabsize=2
}

%--------------------------------------

\title{Longest Common Subsequence}
\author{Шерстобитов Андрей \\
         332 группа}

\begin{document}
\maketitle
    \newpage
\tableofcontents
    \newpage
\section{Aлгоритм в "дискретном"\ случае}

\subsection{Постановка задачи}
 Для данных последовательностей найти самую длинную общую подпоследовательность, при этом
подпоследовательность может быть разрывна в данной последовательности.

Приведу пример:

Пусть даны две последовательности:
\begin{align*}
    S1 &= \{B, C, D, A, A, C, D\} \\
    S2 &= \{A, C, D, B, A, C\}
\end{align*}

Тогда их общие подпоследовательности это
$$\{B, C\},\ \{C, D, A, C\},\ \{D, A, C\},\ \{A, A, C\},\ \{A, C\},\ \{C, D\},\ ...$$
Среди этих подпоследовательностей мы найдем
самую длинную, а именно $\{C, D, A, C\}$
\subsection{Общие понятия}
Прежде чем переходить непосредственно к алгоритму,
    рассмотрим несколько понятий:

\begin{definition}
    Назовем $\Omega$ множество элементов,
    из которых мы будем составлять наши последовательности.
    В данном разделе легче представлять это множество, как множество
    букв английского алфавита, а последовательности - соответственно - строки.
\end{definition}
\begin{definition}
    Знаком $\&$ будем обозначать стандартную конкатенацию последовательностей.
    Например: $\{a,b,c\}\ \&\ \{c,b,a\}\ =\ \{a,b,c,c,b,a\}$.
\end{definition}
\begin{definition}
    Префиксом последовательности $S=\{x_1,\ldots,x_n\}$ длины $0 < k\leq n$
    обозначим $S_k=\{x_1,\ldots, x_k\}$
\end{definition}
\begin{definition}
    Назовем $\LCS(S1, S2)$ функцию, которая принимает на вход
    две последовательности $S1 = \{x_1, \ldots, x_n\}$ и $S2 = \{y_1, \ldots, y_n\}$ и возвращает
    наибольшую общую подпоследовательность $C$.
\end{definition}
    Докажем вспомогательную лемму:
\begin{lemma} \label{lm::props}
    Функция $\LCS(S1, S2)$ обладает следующими свойствами:
    \begin{enumerate}
        \item $\forall\ A\in\Omega$
            $$\LCS(S1\ \&\ A,\ S2\ \&\ A) = \LCS(S1,\ S2)\ \&\ A$$
        \item $\forall A\neq B\in\Omega$
            $$\LCS(S1\ \&\ A, S2\ \&\ B) = \max \{\LCS(S1\ \&\ A, S2), \LCS(S1, S2\ \&\ B)\}$$
    \end{enumerate}
\end{lemma}
\begin{proof}
    \begin{enumerate}
        \item Представим нахождение общей последовательности как
        последовательное сопоставление элементов последовательностей друг за
        другом без пересечений.

        Рассмотрим $S1 = \{a, b, c, d, e, f\},\ S2 = \{d, a, c, e, q\}$ и уже
        найденную для них общую подпоследовательность $C = \{a, c, e\}$, тогда \par
        $$\begin{tabular}{cccccc}
            a\tikzmark{fa1} & b & c\tikzmark{fc1} & d & e \tikzmark{fe1} & f  \\
            d & a\tikzmark{sa1} & c\tikzmark{sc1} & e \tikzmark{se1} & q  \\
          \end{tabular}
          \begin{tikzpicture}[overlay, remember picture, shorten >=7pt, shorten <=4pt, transform canvas={yshift=0.1\baselineskip}]
            \draw [->] ({pic cs:fa1}) -- ({pic cs:sa1});
            \draw [->] ([xshift=-2.5pt]{pic cs:fc1}) -- ([xshift=-2.5pt]{pic cs:sc1});
            \draw [->] ([xshift=-4pt]{pic cs:fe1}) -- ([xshift=-7pt]{pic cs:se1});
          \end{tikzpicture}
        $$
        Добавив в конец каждой последовательности по одному элементу (пусть $p$), мы не нарушим
        наш алгоритм, так как новых пересечений быть не может.
        $$\begin{tabular}{ccccccc}
            a\tikzmark{fa} & b & c\tikzmark{fc} & d & e \tikzmark{fe} & f &p \tikzmark{fp}\\
            d & a\tikzmark{sa} & c\tikzmark{sc} & e \tikzmark{se} & q &p \tikzmark{sp}\\
          \end{tabular}
          \begin{tikzpicture}[overlay, remember picture, shorten >=7pt, shorten <=4pt, transform canvas={yshift=0.1\baselineskip}]
            \draw [->] ({pic cs:fa}) -- ({pic cs:sa});
            \draw [->] ([xshift=-2.5pt]{pic cs:fc}) -- ([xshift=-2.5pt]{pic cs:sc});
            \draw [->] ([xshift=-4pt]{pic cs:fe}) -- ([xshift=-7pt]{pic cs:se});
            \draw [->] ([xshift=-4pt]{pic cs:fp}) -- ([xshift=-7pt]{pic cs:sp});
          \end{tikzpicture}
        $$
        \item В противном случае, если последние элементы последовательностей
        $S1 =\{x_1, \ldots,x_n\}$, $S2=\{y_1,\ldots,y_m\}$ различны $(x_n\neq y_m)$,
        то какой-либо из них не входит в наибольшую \underline{общую}
        подпоследовательность, иначе могут появиться самопересечения. Значит
        нужно рассмотреть последовательности $S1 =\{x_1, \ldots,x_n\}$, $S2'=\{y_1,\ldots,y_{m-1}\}$
        и $S1' =\{x_1, \ldots,x_{n-1}\}$, $S2=\{y_1,\ldots,y_{m}\}$ и среди них
        выбрать максимальную общую подпоследовательность.
    \end{enumerate}
\end{proof}
    Таким образом для последователностей $X = (x_1,\ldots,x_n), Y=(y_1,\ldots,y_m)$,
    их префиксов $X_{1,2,\ldots,n}, Y_{1,2,\ldots,m}$ наша функция $\LCS(X_i,Y_j)$ принимает вид:
    \begin{equation*}
        \LCS(X_i,Y_j) = \begin{cases}
                \emptyset & \text{если } i = 0 \text{ или } j = 0 \\
                \LCS(X_{i-1}, Y_{j-1})\&x_i & \text{если } i,j>0 \text{ и } x_i=y_j \\
                \max\{\LCS(X_{i}, Y_{j-1}), \LCS(X_{i-1}, Y_{j})\} & \text{если } i,j>0 \text{ и } x_i\neq y_j
            \end{cases}
    \end{equation*}

\subsection{Алгоритм}
Рассмотрим алгоритм для последовательностей $X = (x_1,\ldots,x_n)$,
$Y=(y_1,\ldots,y_m)$:
\begin{enumerate}
    \item Создаем таблицу размера $(n+1)\times (m+1)$, где
        первая строка и первый столбец заполнены нулями.
        Это нужно для случая, когда одна из последовательностей пустая.
    \item Заполняем эту таблицу следующим образом:
        \begin{enumerate}
            \item Если элемент на текущей строке и текущей колонке
                совпадают, то по свойству 1 заполняем ячейку числом на 1 больше,
                чем в левой верхней диагональной ячейке. Проводим
                стрелку к этой же ячейке от текущей.
            \item Иначе по свойству 2 берем максимальное значние с предыдущей колонки
                и предыдущей строки и заполняем текущую ячейку. Проводим стрелку
                к максимальному элементу. Если они равны, то к любому из них.
        \end{enumerate}
    \item Повторяем предыдущий шаг до заполнения таблицы.
    \item Получаем в правой нижней ячейке таблицы длину максимальной общей подпоследовательности.
    \item Чтобы распечатать искомую подпоследовательность, начиинаем
        с правого нижнего элемента и двигаемся по стрелкам. Добравшись
        до первой строки или столбца получим нашу последовательность.
\end{enumerate}
\subsection{Примеры}
    Рассмотирм две последовательности $X = \{C, B, D, A\}, Y = \{A,C,A,D,B\}$ и выполним
    для них алгоритм.
    \begin{enumerate}
        \item Составляем требуемую таблицу.
    $$
    \begin{tikzpicture}
        \matrix(m)[table] {
                    & $\emptyset$ & $C$ & $B$ & $D$ & $A$ \\
        $\emptyset$ &       0     &  0  &  0  &  0  & 0 \\
          $A$       &       0     &     &     &     & \\
          $C$       &       0     &     &     &     & \\
          $A$       &       0     &     &     &     & \\
          $D$       &       0     &     &     &     & \\
          $B$       &       0     &     &     &     & \\
        };
    \end{tikzpicture}
    $$
    \item Заполним первую строку соответственно нашему алгоритму и
    свойствам нашей функции LCS(S1, S2).
    $$
    \begin{tikzpicture}
        \matrix(m)[table] {
                    & $\emptyset$ & $C$ & $B$ & $D$ & $A$ \\
        $\emptyset$ &       0     &  0  &  0  &  0  & 0  \\
          $A$       &       0     &  0  &  0  &  0  & 1  \\
          $C$       &       0     &     &     &     & \\
          $A$       &       0     &     &     &     & \\
          $D$       &       0     &     &     &     & \\
          $B$       &       0     &     &     &     & \\
        };
        % the arrows
        \begin{scope}[shorten >=7pt,shorten <= 7pt]
        \draw[->]  (m-3-3.center) -- (m-2-3.center);
        \draw[->]  (m-3-3.center) -- (m-3-2.center);

        \draw[->]  (m-3-4.center) -- (m-3-3.center);
        \draw[->]  (m-3-4.center) -- (m-2-4.center);

        \draw[->]  (m-3-5.center) -- (m-3-4.center);
        \draw[->]  (m-3-5.center) -- (m-2-5.center);

        \draw[->]  (m-3-6.center) -- (m-2-5.center);
        \end{scope}
        \end{tikzpicture}
    $$
        \item Повторим для последующих строк.
    $$
        \begin{tikzpicture}
            % the matrix entries
            \matrix (m) [table] {
                        & $\emptyset$ & $C$ & $B$ & $D$ & $A$ \\
            $\emptyset$ &       0     &  0  &  0  &  0  & 0  \\
              $A$       &       0     &  0  &  0  &  0  & 1  \\
              $C$       &       0     &  1  &  1  &  1  & 1  \\
              $A$       &       0     &  1  &  1  &  1  & 2  \\
              $D$       &       0     &  1  &  1  &  2  & 2  \\
              $B$       &       0     &  1  &  2  &  2  & 2  \\
            };
            % the arrows
            \begin{scope}[shorten >=7pt,shorten <= 7pt]
            \draw[->]  (m-3-3.center) -- (m-2-3.center);
            \draw[->]  (m-3-3.center) -- (m-3-2.center);
            \draw[->]  (m-3-4.center) -- (m-3-3.center);
            \draw[->]  (m-3-4.center) -- (m-2-4.center);
            \draw[->]  (m-3-5.center) -- (m-3-4.center);
            \draw[->]  (m-3-5.center) -- (m-2-5.center);
            \draw[->]  (m-3-6.center) -- (m-2-5.center);

            \draw[->]  (m-4-3.center) -- (m-3-2.center);
            \draw[->]  (m-4-4.center) -- (m-4-3.center);
            \draw[->]  (m-4-5.center) -- (m-4-4.center);
            \draw[->]  (m-4-6.center) -- (m-4-5.center);
            \draw[->]  (m-4-6.center) -- (m-3-6.center);
    
            \draw[->]  (m-5-3.center) -- (m-4-3.center);
            \draw[->]  (m-5-4.center) -- (m-4-4.center);
            \draw[->]  (m-5-4.center) -- (m-5-3.center);
            \draw[->]  (m-5-5.center) -- (m-4-5.center);
            \draw[->]  (m-5-5.center) -- (m-5-4.center);
            \draw[->]  (m-5-6.center) -- (m-4-5.center);

            \draw[->]  (m-6-3.center) -- (m-5-3.center);
            \draw[->]  (m-6-4.center) -- (m-5-4.center);
            \draw[->]  (m-6-4.center) -- (m-6-3.center);
            \draw[->]  (m-6-5.center) -- (m-5-4.center);
            \draw[->]  (m-6-6.center) -- (m-5-6.center);
            \draw[->]  (m-6-6.center) -- (m-6-5.center);

            \draw[->]  (m-7-3.center) -- (m-6-3.center);
            \draw[->]  (m-7-4.center) -- (m-6-3.center);
            \draw[->]  (m-7-5.center) -- (m-6-5.center);
            \draw[->]  (m-7-5.center) -- (m-7-4.center);
            \draw[->]  (m-7-6.center) -- (m-6-6.center);
            \draw[->]  (m-7-6.center) -- (m-7-5.center);
            \end{scope}
        \end{tikzpicture}
    $$
        \item Значение в последней ячейке - длина наибольшей общей последовательности.
    $$
        \begin{tikzpicture}
            % the matrix entries
            \matrix (m) [table] {
                        & $\emptyset$ & $C$ & $B$ & $D$ & $A$ \\
            $\emptyset$ &       0     &  0  &  0  &  0  & 0  \\
              $A$       &       0     &  0  &  0  &  0  & 1  \\
              $C$       &       0     &  1  &  1  &  1  & 1  \\
              $A$       &       0     &  1  &  1  &  1  & 2  \\
              $D$       &       0     &  1  &  1  &  2  & 2  \\
              $B$       &       0     &  1  &  2  &  2  & |[fill=gray!20]|2  \\
            };
            % the arrows
            \begin{scope}[shorten >=7pt,shorten <= 7pt]
            \draw[->]  (m-3-3.center) -- (m-2-3.center);
            \draw[->]  (m-3-3.center) -- (m-3-2.center);
            \draw[->]  (m-3-4.center) -- (m-3-3.center);
            \draw[->]  (m-3-4.center) -- (m-2-4.center);
            \draw[->]  (m-3-5.center) -- (m-3-4.center);
            \draw[->]  (m-3-5.center) -- (m-2-5.center);
            \draw[->]  (m-3-6.center) -- (m-2-5.center);

            \draw[->]  (m-4-3.center) -- (m-3-2.center);
            \draw[->]  (m-4-4.center) -- (m-4-3.center);
            \draw[->]  (m-4-5.center) -- (m-4-4.center);
            \draw[->]  (m-4-6.center) -- (m-4-5.center);
            \draw[->]  (m-4-6.center) -- (m-3-6.center);
    
            \draw[->]  (m-5-3.center) -- (m-4-3.center);
            \draw[->]  (m-5-4.center) -- (m-4-4.center);
            \draw[->]  (m-5-4.center) -- (m-5-3.center);
            \draw[->]  (m-5-5.center) -- (m-4-5.center);
            \draw[->]  (m-5-5.center) -- (m-5-4.center);
            \draw[->]  (m-5-6.center) -- (m-4-5.center);

            \draw[->]  (m-6-3.center) -- (m-5-3.center);
            \draw[->]  (m-6-4.center) -- (m-5-4.center);
            \draw[->]  (m-6-4.center) -- (m-6-3.center);
            \draw[->]  (m-6-5.center) -- (m-5-4.center);
            \draw[->]  (m-6-6.center) -- (m-5-6.center);
            \draw[->]  (m-6-6.center) -- (m-6-5.center);

            \draw[->]  (m-7-3.center) -- (m-6-3.center);
            \draw[->]  (m-7-4.center) -- (m-6-3.center);
            \draw[->]  (m-7-5.center) -- (m-6-5.center);
            \draw[->]  (m-7-5.center) -- (m-7-4.center);
            \draw[->]  (m-7-6.center) -- (m-6-6.center);
            \draw[->]  (m-7-6.center) -- (m-7-5.center);
            \end{scope}
        \end{tikzpicture}
    $$
        \item Находим общие подпоследовательности, двигаясь по стрелочкам.
    $$
        \begin{tikzpicture}
            % the matrix entries
            \matrix (m) [table] {
                        & $\emptyset$ & $C$ & $B$ & $D$ & $A$ \\
            $\emptyset$ &       0     &  0  &  0  &  0  & 0  \\
              $A$       & |[fill=red!20!yellow!20!blue!20!]|0     &  0  &  0  &  0  & 1  \\
              $C$       &       0     &  |[fill=red!20!yellow!20!blue!20!]|1  &  |[fill=yellow!20!blue!20!]|1  &  |[fill=yellow!20]|1  & 1  \\
              $A$       &       0     &  |[fill=red!20!blue!20!]|1  &  |[fill=blue!20!]|1  &  1  & |[fill=yellow!20]|2  \\
              $D$       &       0     &  |[fill=red!20]|1  &  1  &  |[fill=blue!20!]|2  & |[fill=yellow!20!blue!20!]|2  \\
              $B$       &       0     &  1  &  |[fill=red!20]|2  &  |[fill=red!20!blue!20!]|2  & |[fill=red!20!yellow!20!blue!20!]|2  \\
            };
            % the arrows
            \begin{scope}[shorten >=7pt,shorten <= 7pt]
            \draw[->]  (m-3-3.center) -- (m-2-3.center);
            \draw[->]  (m-3-3.center) -- (m-3-2.center);
            \draw[->]  (m-3-4.center) -- (m-3-3.center);
            \draw[->]  (m-3-4.center) -- (m-2-4.center);
            \draw[->]  (m-3-5.center) -- (m-3-4.center);
            \draw[->]  (m-3-5.center) -- (m-2-5.center);
            \draw[->]  (m-3-6.center) -- (m-2-5.center);

            \draw[->]  (m-4-3.center) -- (m-3-2.center);
            \draw[->]  (m-4-4.center) -- (m-4-3.center);
            \draw[->]  (m-4-5.center) -- (m-4-4.center);
            \draw[->]  (m-4-6.center) -- (m-4-5.center);
            \draw[->]  (m-4-6.center) -- (m-3-6.center);
    
            \draw[->]  (m-5-3.center) -- (m-4-3.center);
            \draw[->]  (m-5-4.center) -- (m-4-4.center);
            \draw[->]  (m-5-4.center) -- (m-5-3.center);
            \draw[->]  (m-5-5.center) -- (m-4-5.center);
            \draw[->]  (m-5-5.center) -- (m-5-4.center);
            \draw[->]  (m-5-6.center) -- (m-4-5.center);

            \draw[->]  (m-6-3.center) -- (m-5-3.center);
            \draw[->]  (m-6-4.center) -- (m-5-4.center);
            \draw[->]  (m-6-4.center) -- (m-6-3.center);
            \draw[->]  (m-6-5.center) -- (m-5-4.center);
            \draw[->]  (m-6-6.center) -- (m-5-6.center);
            \draw[->]  (m-6-6.center) -- (m-6-5.center);

            \draw[->]  (m-7-3.center) -- (m-6-3.center);
            \draw[->]  (m-7-4.center) -- (m-6-3.center);
            \draw[->]  (m-7-5.center) -- (m-6-5.center);
            \draw[->]  (m-7-5.center) -- (m-7-4.center);
            \draw[->]  (m-7-6.center) -- (m-6-6.center);
            \draw[->]  (m-7-6.center) -- (m-7-5.center);
            \end{scope}
        \end{tikzpicture}
        $$
    \end{enumerate}
    Таким образом, получившиеся наибольшие общие подпоследовательности:
    $$\{C, B\}, \{C, A\}, \{C, D\}$$
    \begin{remark}
        Данный пример показал, что наибольшая общая подпоследовательность
        не единственна.
    \end{remark}

\newpage
\subsection{Реализация алгоритма на C++}
\begin{lstlisting}
#include <vector>

template<typename T>
using Sequence<T> = std::vector<T>;

template<typename T>
std::vector<T> LCS(const std::vector<T>& S1, const std::vector<T>& S2) {
    std::vector<std::vector<int>> table(S1.size() + 1,
                                    std::vector<int>(S2.size() + 1));
    int i, j;

    for (i = 0; i <= a.size(); ++i) {
        for (j = 0; j <= b.size(); ++j) {
            if (i == 0 || j == 0)
                table[i][j] = 0;
            else if (S1[i - 1] == S2[j - 1])
                table[i][j] = table[i - 1][j - 1] + 1;
            else
                table[i][j] = std::max(table[i - 1][j], table[i][j - 1]);
        }
    }

    int index = table[S1.size()][S2.size()]
    std::vector<T> LCS(index);

    while (i > 0 && j > 0) {
        if (S1[i - 1] == S2[j - 1]) {
          LCS[index - 1] = S1[i - 1];
          i--, j--, index--;
        } else if (LCS_table[i - 1][j] > LCS_table[i][j - 1]) {
            i--;
        } else {
            j--;
        }
    }

    return LCS;
}
\end{lstlisting}

\textbf{Сложность работы:} $\mathcal{O}(mn)$, где $m$ и $n$ - длины последовательностей.

\textbf{Сложнсоть по памяти:} $\mathcal{O}(mn)$, где $m$ и $n$ - длины последовательностей.

\end{document}